\subsection{A Hybrid LP-Separator Recursion Approach}

Assume $0<\rho<\eta<3\rho<1$, and put
\[
    \delta=\eta-\rho,
    \qquad
    x=\ln\left(\frac{2\rho}{\delta}\right),
    \qquad
    k=\left\lfloor x\right\rfloor.
\]
The hybrid first applies anchored-LP rounding with target $\tau=\rho+e^k\delta$, then runs separator recursion on the resulting components. Since $e^k\delta\leq2\rho$, the LP target satisfies $\tau\leq3\rho<1$, and at most $k$ separator rounds suffice.

\subsubsection{The hybrid procedure}
\begin{algorithm}[tbp]
    \DontPrintSemicolon
    \caption{Hybrid LP-separator recursion}\label{alg:hybrid-lp-centroid}
    \textbf{Input} Supply tree $T=(V,E,c)$, demand graph $(V,D,w)$, and fixed rational parameters $0<\rho<\eta<3\rho<1$\;
    \textbf{Output} An $\eta$-feasible supply-edge set $F\subseteq E$\;
    \If{$W=0$}{
        \textbf{return} $F=\emptyset$\;
    }
    Set $\delta\gets\eta-\rho$, $x\gets\ln(2\rho/\delta)$, and $k\gets\lfloor x\rfloor$\;
    Set $\tau\gets\rho+e^k\delta$\;
    Run the anchored-LP algorithm with feasibility target $\tau$. Let $F_0$ be the returned cut and let $\cH_0$ be the components of $T\setminus F_0$\;
    Set $B_\rho\gets\rho W$, $B_\eta\gets\eta W$, and $F\gets F_0$\;
    \ForEach{$H\in\cH_0$ with $w\br{V_H}>B_\eta$}{
        Set $\rho_H\gets B_\rho/w\br{V_H}$ and $\eta_H\gets B_\eta/w\br{V_H}$\;
        Run Algorithm~\ref{alg:centroid-max-coverage} on the instance induced by $V_H$ with parameters $\rho_H,\eta_H$, obtaining a supply-edge set $F_H$\;
        $F\gets F\cup F_H$\;
    }
    \textbf{return} $F$\;
\end{algorithm}

Apply Corollary~\ref{cor:tree-general-demand-bicriteria} with feasibility target $\tau=\rho+e^k\delta$, and let $F_0$ be the resulting supply-edge set and $\cH_0$ the components of $T\setminus F_0$. Every $H\in\cH_0$ has demand at most $\tau W$, and the normalization cost satisfies
\begin{equation}\label{eq:tree-hybrid-normalization-cost}
    c\br{F_0}
    \leq
    \left(2+\frac{2\rho}{e^k\delta}\right)\cdot\OPT_{B_\rho}(T).
\end{equation}

For each $H\in\cH_0$ with $w\br{V_H}>B_\eta$, run Algorithm~\ref{alg:centroid-max-coverage} on the instance induced by $V_H$ with local parameters
\[
    \rho_H=\frac{B_\rho}{w\br{V_H}},
    \qquad
    \eta_H=\frac{B_\eta}{w\br{V_H}}.
\]
These parameters make the algorithm's local budgets equal to $B_\rho$ and $B_\eta$, respectively. Let $F_H$ be its returned supply-edge set, and set $F=F_0\cup\bigcup_{H\in\cH_0}F_H$.

\begin{theorem}\label{thm:tree-hybrid-recursion}
    Fix rational parameters $0<\rho<\eta<3\rho<1$. For integral demand weights on a supply tree, the hybrid algorithm returns an $\eta$-feasible supply-edge set $F$ satisfying
    \[
        c\br{F}\leq\Gamma_{\mathrm{H}}(\rho,\eta)\cdot\OPTdem{T}{\rho},
    \]
    where
    \begin{equation}\label{eq:tree-hybrid-ratio}
        \Gamma_{\mathrm{H}}(\rho,\eta)
        =e+1+\ln\left(\frac{2\rho}{\eta-\rho}\right).
    \end{equation}
    For fixed $\rho$ and $\eta$, the algorithm runs in polynomial time.
\end{theorem}
\begin{proof}
    Let $\delta=\eta-\rho$, $x=\ln(2\rho/\delta)$, $k=\lfloor x\rfloor$, and $s=x-k\in[0,1)$. Since $k\leq x$, we have $e^k\delta\leq2\rho$ and hence $\tau=\rho+e^k\delta\leq3\rho<1$. By Corollary~\ref{cor:tree-general-demand-bicriteria}, the normalization cut satisfies~\eqref{eq:tree-hybrid-normalization-cost}.

    By the LP normalization, $w\br{V_H}\leq\tau W=(\rho+e^k\delta)W$, so every normalized component's excess above $B_\rho$ is at most $e^k\delta W$. Each separator-coverage round reduces the excess of every child by a factor of at least $e$. Hence after at most $k$ rounds all components have excess at most $\delta W=B_\eta-B_\rho$ and are $\eta$-feasible. At each depth, the processed components are pairwise edge-disjoint, and restricting an optimal $B_\rho$-feasible cut of $T$ to them shows that the total cost charged at that depth is at most $\OPT_{B_\rho}(T)$. The separator recursion therefore adds cost at most $k\OPT_{B_\rho}(T)$.

    Combining this with~\eqref{eq:tree-hybrid-normalization-cost} gives
    \[
        c\br{F}\leq
        \left(k+2+\frac{2\rho}{e^k\delta}\right)\OPT_{B_\rho}(T).
    \]
    Since $x=k+s$ and $e^k\delta=2\rho/e^s$, the coefficient is
    \[
        k+2+\frac{2\rho}{e^k\delta}
        =k+2+e^s
        =x+2+e^s-s
        \leq x+e+1,
    \]
    where $e^s-s\leq e-1$ for $0\leq s<1$. This proves the claimed bound. Since all of the constants in this construction ultimately depend only on the fixed parameters $\rho$ and $\eta$, the algorithm runs in polynomial time.
\end{proof}